\documentclass[11pt]{article}
\usepackage[margin=1in]{geometry}
\usepackage{amsmath,amssymb,amsthm}
\usepackage{hyperref}
\usepackage{enumitem}

\newtheorem{theorem}{Theorem}
\newtheorem{lemma}[theorem]{Lemma}
\newtheorem{corollary}[theorem]{Corollary}
\theoremstyle{definition}
\newtheorem{definition}[theorem]{Definition}
\theoremstyle{remark}
\newtheorem{remark}[theorem]{Remark}

\usepackage{xurl}
\let\oldus\_
\renewcommand{\_}{\oldus\allowbreak}

\title{A proof of Conjecture 00000000206\\[2pt]
\large Every value occurs only finitely often in Recam\'an's sequence}
\date{}

\begin{document}
\sloppy
\maketitle

\section{The conjecture}

Conjecture 00000000206 reads:

\begin{quote}
\emph{Definition:} The Recam\'an sequence is defined by the backtracking-doubling recurrence.
\emph{Conjecture:} Every nonnegative integer occurs only finitely often in the sequence; and the
finiteness is settled by interval covering.
\end{quote}

\paragraph{Reading.} The Recam\'an sequence is OEIS A005132. The phrase ``backtracking-doubling'' is
loose wording: no doubling occurs in the standard sequence, and we use the standard definition below.
The first clause is the mathematical claim: for every integer $m \ge 0$, the set $\{n : a_n = m\}$ is
finite. The second clause, ``the finiteness is settled by interval covering'' (Chinese text: the core of
the finiteness is interval covering), names no object or property. We read it as a remark on the method
of proof, not as a mathematical statement, and we do not formalize it. The proof below is an interval
argument: every add-step occurrence of $m$ lies at an index in $[1, m]$.

\section{Definition}

The retrieved Wikipedia article on Recam\'an's sequence
(\url{https://en.wikipedia.org/wiki/Recam%C3%A1n%27s_sequence}) defines $a_0 = 0$ and, for $n \ge 1$,
\[
a_n = \begin{cases} a_{n-1} - n & \text{if } a_{n-1} - n > 0 \text{ and is not already in the sequence},\\
a_{n-1} + n & \text{otherwise.}\end{cases}
\]
OEIS A005132 (\url{https://oeis.org/A005132}) states the same rule with ``nonnegative'' in place of
``$> 0$''. ``Not already in the sequence'' means not among $a_0, \dots, a_{n-1}$. The two guards give the
same sequence, because the value $0 = a_0$ is always already present.

In Lean (ASCII rendering; indices shifted to $n+1$, values in $\mathbb{N}$, so the subtraction is the true difference
under the guard):
\begin{verbatim}
def IsRecaman (a : Nat -> Nat) : Prop :=
  a 0 = 0 /\ forall n, a (n + 1) =
    if n + 1 < a n /\ forall k <= n, a k != a n - (n + 1)
    then a n - (n + 1) else a n + (n + 1)
\end{verbatim}
\texttt{IsRecamanOEIS} is the same with \texttt{n + 1 <= a n}. The sequence \texttt{recaman} is built by
list recursion (\texttt{recList n} $= [a_n, \dots, a_0]$). We prove \texttt{isRecaman\_recaman},
\texttt{isRecamanOEIS\_recaman} (it satisfies both forms) and \texttt{isRecaman\_unique} (the recurrence
determines the sequence). The first 27 terms are checked by \texttt{decide} against the Wikipedia list
$0, 1, 3, 6, 2, 7, 13, 20, 12, 21, 11, 22, 10, 23, 9, 24, 8, 25, 43, 62, 42, 63, 41, 18, 42, 17, 43$;
note $a_{20} = a_{24} = 42$, so values can repeat.

\section{Proof}

\begin{lemma}[\texttt{finite\_of\_step}]
Let $a : \mathbb{N} \to \mathbb{N}$ be such that for every $n$, either $a_{n+1} = a_n + (n+1)$ or
$a_{n+1} \ne a_k$ for all $k \le n$. Then for every $m$ the set $\{n : a_n = m\}$ is finite.
\end{lemma}

\begin{proof}
Let $n > m$ with $a_n = m$; then $n = p + 1$. If $a_n = a_p + n$ then $m \ge n$, a contradiction. So
$a_n \ne a_k$ for all $k < n$, i.e.\ $n$ is the first index with value $m$. Hence at most one index
$N > m$ has $a_N = m$. If there is one, $\{n : a_n = m\} \subseteq [0, N]$; otherwise $\{n : a_n = m\}
\subseteq [0, m]$.
\end{proof}

\begin{theorem}[\texttt{finite\_occurrences\_of\_isRecaman}, \texttt{finite\_occurrences\_of\_isRecamanOEIS},
\texttt{recaman\_finite\_occurrences}]
For every sequence satisfying either form of the recurrence, and in particular for \texttt{recaman}, every
value $m \in \mathbb{N}$ occurs only finitely often.
\end{theorem}

\begin{proof}
In the subtract branch the guard says $a_{n+1} = a_n - (n+1)$ is not among $a_0, \dots, a_n$; in the
other branch $a_{n+1} = a_n + (n+1)$. So the lemma applies.
\end{proof}

\section{Lean formalization}

File \texttt{lean/Conjecture206/Basic.lean} (namespace \texttt{Conjecture206}, 166 lines, only import
\texttt{Mathlib}). Main theorem (ASCII rendering):
\begin{verbatim}
theorem recaman_finite_occurrences (m : Nat) : {n | recaman n = m}.Finite
\end{verbatim}
Supporting results: \texttt{finite\_of\_step}, \texttt{finite\_occurrences\_of\_isRecaman},
\texttt{finite\_occurrences\_of\_isRecamanOEIS}, \texttt{recList\_succ}, \texttt{mem\_recList},
\texttt{isRecaman\_recaman}, \texttt{isRecamanOEIS\_recaman}, \texttt{isRecaman\_unique}.

\paragraph{Axiom audit.} \texttt{\#print axioms} for \texttt{recaman\_finite\_occurrences},
\texttt{finite\_occurrences\_of\_isRecamanOEIS} and \texttt{isRecamanOEIS\_recaman} reports only
\texttt{propext}, \texttt{Classical.choice}, \texttt{Quot.sound} (\texttt{isRecaman\_unique}: only
\texttt{propext}, \texttt{Quot.sound}). There is no \texttt{sorry} and no \texttt{native\_decide}.

\end{document}
